% Authoritative exercise source. Preserve all derivations and diagrams.

\ExerciseStart{G3.03}{Applicare X, Z e H agli stessi stati}{Guida esplicita}

\BookHeading{Problema, dati e obiettivo}

Su tre copie della stessa preparazione \(\ket\psi=(\sqrt3,1)^T/2\), applicare separatamente X, Z e H. Calcolare ampiezze e probabilità, quindi confrontare la geometria delle tre trasformazioni.

\ExTheory{\ExRef{c3-porte}{3.03}}{\(X=\begin{pmatrix}0&1\\1&0\end{pmatrix}\), \(Z=\begin{pmatrix}1&0\\0&-1\end{pmatrix}\), \(H=\frac1{\sqrt2}\begin{pmatrix}1&1\\1&-1\end{pmatrix}\). I confronti riguardano preparazioni ripetute, non una sequenza X–Z–H.}

\ExPrediction{X scambia le componenti, Z ne cambia un segno, H forma somma e differenza. Nella sezione reale delle ampiezze questi effetti sono visibili come simmetrie rispetto a rette diverse.}

\ExFigure{\begin{tikzpicture}[>=Stealth,exlabel]
\draw[exaxis] (-0.35,0)--(1.55,0) node[right] {$x_1$};\draw[exaxis] (0,-0.85)--(0,1.416025) node[above] {$x_2$};\draw[->,Ink,line width=1.1pt] (0,0)--(0.866025,0.5) node[anchor=south west] {$ \psi $};\draw[->,Accent,line width=1.1pt] (0,0)--(0.5,0.866025) node[anchor=south west] {$ X\psi $};\draw[->,Secondary,line width=1.1pt] (0,0)--(0.866025,-0.5) node[anchor=north west] {$ Z\psi $};
\end{tikzpicture}}{Piano reale delle ampiezze normalizzate: scambio e inversione di segno.}


\begin{ExTasks}

\item Calcolare le probabilità della preparazione.

\item Svolgere entrambe le righe di X.

\item Svolgere entrambe le righe di Z.

\item Svolgere entrambe le righe di H.

\item Controllare le tre norme e confrontare le distribuzioni.

\end{ExTasks}

\ExSolution

\ExStep{Preparazione comune}

\(\|\psi\|^2=(3+1)/4=1\), con \(p=(3/4,1/4)\). Le tre copie partono dagli stessi dati.

\ExStep{Applicare X}

\ExEq{X\ket\psi&=\tfrac12\colv{0\sqrt3+1\cdot1\\1\sqrt3+0\cdot1}
=\tfrac12\colv{1\\\sqrt3}.}
Le probabilità diventano \((1/4,3/4)\).

\ExStep{Applicare Z}

\ExEq{Z\ket\psi&=\tfrac12\colv{1\sqrt3+0\cdot1\\0\sqrt3-1\cdot1}
=\tfrac12\colv{\sqrt3\\-1}.}
I moduli quadri restano \((3/4,1/4)\); la fase relativa cambia di \(\pi\).

\ExStep{Applicare H}

\ExEq{H\ket\psi&=\frac1{2\sqrt2}\colv{\sqrt3+1\\\sqrt3-1},\\p_H(0)&=\frac{3+2\sqrt3+1}{8}=\frac{2+\sqrt3}{4},\\p_H(1)&=\frac{3-2\sqrt3+1}{8}=\frac{2-\sqrt3}{4}.}

\ExStep{Verificare e confrontare}

Per X e Z le somme sono \(1/4+3/4=1\). Per H, \([(2+\sqrt3)+(2-\sqrt3)]/4=1\). La probabilità di zero è rispettivamente \(0{,}25\), \(0{,}75\) e circa \(0{,}9330\). Anche quando Z lascia invariata questa lettura, una successiva H può rivelare il segno cambiato.

\par\noindent\begin{minipage}{\linewidth}
\ExFigure{\begin{tikzpicture}[>=Stealth,exlabel]
\draw[exaxis] (-0.35,0)--(1.55,0) node[right] {$x_1$};\draw[exaxis] (0,-0.35)--(0,1.05) node[above] {$x_2$};\draw[->,Secondary,line width=1.1pt] (0,0)--(0.866025,0.5) node[anchor=south west] {$ \psi $};\draw[->,Accent,line width=1.1pt] (0,0)--(0.965926,0.258819) node[anchor=south west] {$ H\psi $};
\end{tikzpicture}}{La Hadamard conserva la lunghezza e cambia le componenti mediante interferenza.}


\ExInterpret{Le tre matrici preservano la norma ma agiscono diversamente sulle ampiezze. Le rappresentazioni nel piano reale sono valide per questi stati reali; lo spazio completo del qubit comprende fasi complesse.}

\ExLab{Nella sezione 3.03 confrontare la definizione delle porte con gli stati intermedi. Nello strumento di un qubit preparare angolo polare \(60^\circ\), fase zero, e ripristinare la preparazione prima di ogni porta: il confronto richiede tre copie equivalenti.}


\end{minipage}
